% Page 051

\seite{51}

\abschnitt{Spaziergang}
\pstart
Am 3 September machte hiesige Schule einen Spaziergang nach St.\ Thomas

\pend
\abschnitt{Geschehen 8. 1. 14}
\pstart


\pend
\abschnitt{Kirchlichkeiten}
\pstart

Der Geburtstag Sr. Majestät des Kaisers wurde im Beisein ge\ohb
feiert von den 2 Vätern derselben im Saale des Gasthauses. Den\ob
Kindern wurde{[n]} Geschenke\unsicher{} es wurden vaterländische Lieder gesungen\ob
und die Feste vorgetragen. Die Aufsätze an der Kinder galt\ob
jener Krone von Kuchen.

\pend
\abschnitt{Osterprüfung}
\pstart

wurde am 8.\ März durch den Ortsschulinspek{[tor]} gegen Abend\unsicher\ob
mit Segen vorgehalten.

\pend
\jahresueberschrift{Das Schuljahr 1914/15}
\pstart

Aufgelösen wurden in hiesigen Schule 4 Kinder 2 Knaben\ob
2 Mädchen. Darüber mit denen für die Landwirtschaft\ob
die Genommen haben an 9 Kinder 2 stärkere Klasse.\ob
Die Schülerzahl beträgt 47, 11 Knaben, 11 Mädchen.\ob
Alle Kinder sind katholisch.


Besondere Ereignisse\ob
Die zweite Hälfte des Jahres 1914 brachte Deutschland große Nöte, der\ob
Krieg brach auf das Volk es galten den Franzosen und dessen\unsicher{} Ge\ohb
mahlen zu ermahnet. Da die Reden die Bedingungen nicht erfüllten, die\ob
Slaverei von ihnen forderte, daß des König in Deutschland, als Vater\ob
leider Österreichs, nur bereit, zu stehen. Die Stimmen der Russen,\ob
England, Frankreich hatten Längs einen Kniff und Deutschlander\ob
die gerne in aller gebildeter Weise dem Lande den Frieden erhalten\ob
es war außen eifrig. Den erlosch ja, das Mounts zu großen hat\ob
und suchte einst Kaiser und Deutschland, als Rußland bereit das\ob
Volk gut der Waffen gerufen hatte. Es waren gen Tage das\ob
Volk mit Angst, alle immer sichter nun, es warte dem Kaiser
\pend
